\chapter{Fazit und Ausblick}
\label{kap:fazit}
\verfasser{\autorSE}

In dieser Arbeit wurde eine webbasierte Spiellobby mit React im Frontend und Node.js mit Express und der Bibliothek \texttt{ws} im Backend entwickelt, um die vier Anforderungen der Aufgabenstellung mit WebSockets umzusetzen.

Echtzeit-Updates sind über die Nachrichten \texttt{lobby\_list} und \texttt{lobby\_updated} umgesetzt, die der Server nach jeder Änderung an alle betroffenen Clients verteilt. Live-Benachrichtigungen erhalten Spieler bei Freundschaftsanfragen, Einladungen, Änderungen des Online-Status und privaten Nachrichten. Das kollaborative Bearbeiten zeigt sich im gemeinsamen Lobby-Chat, im privaten Chat mit Tipp-Anzeige und im gemeinsamen Bereit-Status. Alle Aktualisierungen laufen nach der Anmeldung über eine einzige WebSocket-Verbindung, ohne dass die Seite neu geladen wird. Die Nachweise mit Screenshots, dem Mitschnitt der Nachrichten und den Entwicklerwerkzeugen des Browsers bestätigen dieses Verhalten.

Auf der Serverseite ermöglicht das Single-Thread-Modell von Node.js einen einfachen, sperrfreien Umgang mit dem Zustand im Speicher. Datenbankzugriffe und das Hashen von Passwörtern laufen in Worker Threads, damit sie die Event Loop nicht blockieren. Die Authentifizierung erfolgt nach dem empfohlenen Verfahren für Browser-Sitzungen. Ein kurzlebiges Access-Token und ein rotiertes, serverseitig gespeichertes Refresh-Token liegen in \texttt{HttpOnly}-Cookies mit \texttt{SameSite=Strict}. Der Server prüft Access-Token und Ursprung beim WebSocket-Handshake und bindet die Identität danach fest an die Verbindung. Sitzungen lassen sich damit widerrufen, und aktive Nutzer bleiben trotzdem angemeldet. Auf der Clientseite trennt die Kombination aus einem Hook für die Verbindung und einem externen Zustand-Store die ereignisgetriebene Kommunikation vom deklarativen Rendering.

Für eine Weiterentwicklung bieten sich folgende Punkte an: eine automatische Wiederverbindung mit erneutem Abgleich des Zustands, Ping- und Pong-Frames zum Erkennen abgebrochener Verbindungen, eine Schema-Validierung eingehender Nachrichten, die Übertragung von Änderungen anstelle vollständiger Listen sowie ein gemeinsamer Zustand über mehrere Server-Instanzen. Außerdem könnten Abläufe aus mehreren Datenbankaufrufen, etwa das Senden einer Freundschaftsanfrage, als einzelne Transaktion im Datenbank-Worker ausgeführt werden. Für das gleichzeitige Bearbeiten von Texten könnten \glspl{crdt} ergänzt werden.
